% Options for packages loaded elsewhere
\PassOptionsToPackage{unicode}{hyperref}
\PassOptionsToPackage{hyphens}{url}
\PassOptionsToPackage{dvipsnames,svgnames,x11names}{xcolor}
%
\documentclass[
  letterpaper,
  DIV=11,
  numbers=noendperiod]{scrreprt}

\usepackage{amsmath,amssymb}
\usepackage{iftex}
\ifPDFTeX
  \usepackage[T1]{fontenc}
  \usepackage[utf8]{inputenc}
  \usepackage{textcomp} % provide euro and other symbols
\else % if luatex or xetex
  \usepackage{unicode-math}
  \defaultfontfeatures{Scale=MatchLowercase}
  \defaultfontfeatures[\rmfamily]{Ligatures=TeX,Scale=1}
\fi
\usepackage{lmodern}
\ifPDFTeX\else  
    % xetex/luatex font selection
\fi
% Use upquote if available, for straight quotes in verbatim environments
\IfFileExists{upquote.sty}{\usepackage{upquote}}{}
\IfFileExists{microtype.sty}{% use microtype if available
  \usepackage[]{microtype}
  \UseMicrotypeSet[protrusion]{basicmath} % disable protrusion for tt fonts
}{}
\makeatletter
\@ifundefined{KOMAClassName}{% if non-KOMA class
  \IfFileExists{parskip.sty}{%
    \usepackage{parskip}
  }{% else
    \setlength{\parindent}{0pt}
    \setlength{\parskip}{6pt plus 2pt minus 1pt}}
}{% if KOMA class
  \KOMAoptions{parskip=half}}
\makeatother
\usepackage{xcolor}
\setlength{\emergencystretch}{3em} % prevent overfull lines
\setcounter{secnumdepth}{5}
% Make \paragraph and \subparagraph free-standing
\makeatletter
\ifx\paragraph\undefined\else
  \let\oldparagraph\paragraph
  \renewcommand{\paragraph}{
    \@ifstar
      \xxxParagraphStar
      \xxxParagraphNoStar
  }
  \newcommand{\xxxParagraphStar}[1]{\oldparagraph*{#1}\mbox{}}
  \newcommand{\xxxParagraphNoStar}[1]{\oldparagraph{#1}\mbox{}}
\fi
\ifx\subparagraph\undefined\else
  \let\oldsubparagraph\subparagraph
  \renewcommand{\subparagraph}{
    \@ifstar
      \xxxSubParagraphStar
      \xxxSubParagraphNoStar
  }
  \newcommand{\xxxSubParagraphStar}[1]{\oldsubparagraph*{#1}\mbox{}}
  \newcommand{\xxxSubParagraphNoStar}[1]{\oldsubparagraph{#1}\mbox{}}
\fi
\makeatother


\providecommand{\tightlist}{%
  \setlength{\itemsep}{0pt}\setlength{\parskip}{0pt}}\usepackage{longtable,booktabs,array}
\usepackage{calc} % for calculating minipage widths
% Correct order of tables after \paragraph or \subparagraph
\usepackage{etoolbox}
\makeatletter
\patchcmd\longtable{\par}{\if@noskipsec\mbox{}\fi\par}{}{}
\makeatother
% Allow footnotes in longtable head/foot
\IfFileExists{footnotehyper.sty}{\usepackage{footnotehyper}}{\usepackage{footnote}}
\makesavenoteenv{longtable}
\usepackage{graphicx}
\makeatletter
\def\maxwidth{\ifdim\Gin@nat@width>\linewidth\linewidth\else\Gin@nat@width\fi}
\def\maxheight{\ifdim\Gin@nat@height>\textheight\textheight\else\Gin@nat@height\fi}
\makeatother
% Scale images if necessary, so that they will not overflow the page
% margins by default, and it is still possible to overwrite the defaults
% using explicit options in \includegraphics[width, height, ...]{}
\setkeys{Gin}{width=\maxwidth,height=\maxheight,keepaspectratio}
% Set default figure placement to htbp
\makeatletter
\def\fps@figure{htbp}
\makeatother
% definitions for citeproc citations
\NewDocumentCommand\citeproctext{}{}
\NewDocumentCommand\citeproc{mm}{%
  \begingroup\def\citeproctext{#2}\cite{#1}\endgroup}
\makeatletter
 % allow citations to break across lines
 \let\@cite@ofmt\@firstofone
 % avoid brackets around text for \cite:
 \def\@biblabel#1{}
 \def\@cite#1#2{{#1\if@tempswa , #2\fi}}
\makeatother
\newlength{\cslhangindent}
\setlength{\cslhangindent}{1.5em}
\newlength{\csllabelwidth}
\setlength{\csllabelwidth}{3em}
\newenvironment{CSLReferences}[2] % #1 hanging-indent, #2 entry-spacing
 {\begin{list}{}{%
  \setlength{\itemindent}{0pt}
  \setlength{\leftmargin}{0pt}
  \setlength{\parsep}{0pt}
  % turn on hanging indent if param 1 is 1
  \ifodd #1
   \setlength{\leftmargin}{\cslhangindent}
   \setlength{\itemindent}{-1\cslhangindent}
  \fi
  % set entry spacing
  \setlength{\itemsep}{#2\baselineskip}}}
 {\end{list}}
\usepackage{calc}
\newcommand{\CSLBlock}[1]{\hfill\break\parbox[t]{\linewidth}{\strut\ignorespaces#1\strut}}
\newcommand{\CSLLeftMargin}[1]{\parbox[t]{\csllabelwidth}{\strut#1\strut}}
\newcommand{\CSLRightInline}[1]{\parbox[t]{\linewidth - \csllabelwidth}{\strut#1\strut}}
\newcommand{\CSLIndent}[1]{\hspace{\cslhangindent}#1}

\KOMAoption{captions}{tableheading}
\makeatletter
\@ifpackageloaded{bookmark}{}{\usepackage{bookmark}}
\makeatother
\makeatletter
\@ifpackageloaded{caption}{}{\usepackage{caption}}
\AtBeginDocument{%
\ifdefined\contentsname
  \renewcommand*\contentsname{Table of contents}
\else
  \newcommand\contentsname{Table of contents}
\fi
\ifdefined\listfigurename
  \renewcommand*\listfigurename{List of Figures}
\else
  \newcommand\listfigurename{List of Figures}
\fi
\ifdefined\listtablename
  \renewcommand*\listtablename{List of Tables}
\else
  \newcommand\listtablename{List of Tables}
\fi
\ifdefined\figurename
  \renewcommand*\figurename{Figure}
\else
  \newcommand\figurename{Figure}
\fi
\ifdefined\tablename
  \renewcommand*\tablename{Table}
\else
  \newcommand\tablename{Table}
\fi
}
\@ifpackageloaded{float}{}{\usepackage{float}}
\floatstyle{ruled}
\@ifundefined{c@chapter}{\newfloat{codelisting}{h}{lop}}{\newfloat{codelisting}{h}{lop}[chapter]}
\floatname{codelisting}{Listing}
\newcommand*\listoflistings{\listof{codelisting}{List of Listings}}
\makeatother
\makeatletter
\makeatother
\makeatletter
\@ifpackageloaded{caption}{}{\usepackage{caption}}
\@ifpackageloaded{subcaption}{}{\usepackage{subcaption}}
\makeatother

\ifLuaTeX
  \usepackage{selnolig}  % disable illegal ligatures
\fi
\usepackage{bookmark}

\IfFileExists{xurl.sty}{\usepackage{xurl}}{} % add URL line breaks if available
\urlstyle{same} % disable monospaced font for URLs
\hypersetup{
  pdftitle={Análise do Cumprimento das Metas de Redução de Mortes no Trânsito dos Municípios Brasileiros até 2023.},
  pdfauthor={Ana Beatriz da S. Marques; Dr.~Jorge Tiago Bastos},
  colorlinks=true,
  linkcolor={blue},
  filecolor={Maroon},
  citecolor={Blue},
  urlcolor={Blue},
  pdfcreator={LaTeX via pandoc}}


\title{Análise do Cumprimento das Metas de Redução de Mortes no Trânsito
dos Municípios Brasileiros até 2023.}
\author{Ana Beatriz da S. Marques \and Dr.~Jorge Tiago Bastos}
\date{2025-05-28}

\begin{document}
\maketitle

\renewcommand*\contentsname{Table of contents}
{
\hypersetup{linkcolor=}
\setcounter{tocdepth}{2}
\tableofcontents
}

\bookmarksetup{startatroot}

\chapter*{Resumo}\label{resumo}
\addcontentsline{toc}{chapter}{Resumo}

\markboth{Resumo}{Resumo}

\begin{center}\rule{0.5\linewidth}{0.5pt}\end{center}

O Plano Nacional de Redução de Mortes e Lesões no Trânsito (PNATRANS),
criado pela Lei Federal nº 13.614/2018, tem como objetivo estabelecer
metas de redução da mortalidade no trânsito para os estados e para o
país entre 2019 e 2028. Em 2021, o plano foi revisado, prorrogando o
prazo das metas até 2030 e reformulando seus pilares de atuação. As
metas de redução foram definidas para os municípios brasileiros através
das técnicas de \emph{clusterização} e \emph{benchmarking}. O objetivo,
agora, é observar o desempenho de cada município até 2023, atráves da
visualização de dados. Para isso, foi criado um Dashboard que contém
todas as informações necessárias para essa análise. O DashBoard para
visualização do desempenho dos municípios está disponível neste link:
\href{https://beatrizmarques.shinyapps.io/cumprimento_metas/}{Dashboard}.

\bookmarksetup{startatroot}

\chapter*{Highlights}\label{highlights}
\addcontentsline{toc}{chapter}{Highlights}

\markboth{Highlights}{Highlights}

\begin{center}\rule{0.5\linewidth}{0.5pt}\end{center}

\begin{itemize}
\tightlist
\item
  35\% dos municípios brasileiros já atingiram a meta estabelecida;
\item
  O estado do Acre obteve o melhor desempenho, com uma redução de 15\%
  das mortes no trânsito e 50\% dos seus municípios com a meta já
  atingida;
\item
  Das 27 unidades da federação, 22 obteviveram um aumento no número de
  mortes no trânsito em 2023;
\item
  Alguns estados apresentaram aumento nas mortes de até 26,45\%,
  indicando um agravamento na situação de segurança viária;
\item
  São Paulo foi a capital com melhor desempenho, registrando uma redução
  de 52\% de mortes no trânsito;
\item
  Em 2023, o Brasil registrou 34599 mortes no trânsito, apontando um
  aumento de 6,66\% em relação à 2020;
\end{itemize}

\bookmarksetup{startatroot}

\chapter*{Sobre o Observatório}\label{sobre-o-observatuxf3rio}
\addcontentsline{toc}{chapter}{Sobre o Observatório}

\markboth{Sobre o Observatório}{Sobre o Observatório}

\begin{center}\rule{0.5\linewidth}{0.5pt}\end{center}

O Observatório Nacional de Segurança Viária é uma instituição social sem
fins lucrativos, dedicada a desenvolver ações que contribuam
efetivamente para a redução dos elevados índices de ocorrências no
trânsito brasileiro. Com esse objetivo, um grupo de profissionais
multidisciplinares decidiu reunir todo o seu conhecimento, experiência e
motivação em um único projeto grandioso e desafiador: mobilizar a
sociedade em prol de um trânsito mais seguro.

\bookmarksetup{startatroot}

\chapter*{Como citar}\label{como-citar}
\addcontentsline{toc}{chapter}{Como citar}

\markboth{Como citar}{Como citar}

\begin{center}\rule{0.5\linewidth}{0.5pt}\end{center}

ONSV; UFPR (2025). Análise do Cumprimento das Metas de Redução de Mortes
no Trânsito dos Municípios Brasileiros até 2023. Observatório Nacional
de Segurança Viária e Universidade Federal do Paraná. Disponível em:
https://onsv.github.io/relatorio-cumprimento-metas/

\begin{figure}

\begin{minipage}{0.36\linewidth}
\begin{center}
\includegraphics{img/onsv-logo.jpg}
\end{center}
\end{minipage}%
%
\begin{minipage}{0.27\linewidth}
\begin{center}
\includegraphics{img/marca_UFPR.png}
\end{center}
\end{minipage}%

\end{figure}%

```

\bookmarksetup{startatroot}

\chapter{Introdução}\label{introduuxe7uxe3o}

O Plano Nacional de Redução de Mortes e Lesões no Trânsito (PNATRANS),
instituído pela Lei nº 13.614/2018, tem como objetivo central reduzir em
pelo menos 50\% o número de mortes no trânsito brasileiro ao longo de
uma década. Após dificuldades iniciais de implementação, o plano foi
revisado em 2021, com participação de mais de 100 especialistas e
representantes de diversas instituições públicas e privadas, resultando
em uma versão mais moderna e aplicável, estruturada em seis pilares
temáticos. A nova versão do PNATRANS incorporou as abordagens de
``Sistemas Seguros'' e ``Visão Zero'', adaptadas à realidade brasileira,
e passou a enfatizar o papel estratégico dos municípios na promoção de
ações locais de segurança viária.

Dentre os avanços trazidos pela revisão, destaca-se a elaboração de
indicadores de desempenho e a proposição de metas específicas de redução
de mortes no trânsito para cada município, considerando as
particularidades regionais. Esses elementos viabilizam a atuação baseada
em evidências e permitem um acompanhamento mais preciso do progresso em
diferentes territórios.

Nesse contexto, este novo estudo tem como objetivo realizar o
monitoramento das metas de redução de mortes no trânsito com base nos
dados mais recentes, referentes ao ano de 2023. A proposta é avaliar o
desempenho dos municípios frente às metas estabelecidas, identificar
quais localidades avançaram significativamente na redução da mortalidade
viária, quais ainda enfrentam desafios relevantes e, assim, contribuir
para o fortalecimento da gestão local da segurança no trânsito. Esse
acompanhamento contínuo é essencial para garantir que os esforços
empreendidos estejam alinhados com os objetivos do PNATRANS e para
promover os ajustes necessários na direção de um trânsito mais seguro
para todos.

\bookmarksetup{startatroot}

\chapter{Metodologia}\label{metodologia}

Este estudo tem como objetivo monitorar o cumprimento das metas de
redução de mortes no trânsito estabelecidas pelo Plano Nacional de
Redução de Mortes e Lesões no Trânsito (PNATRANS), com base nos dados
atualizados para o ano de 2023. A análise foi realizada sobre os
municípios brasileiros que apresentaram registros consistentes de óbitos
no trânsito no período base.

A unidade de análise adotada foram os municípios brasileiros que
possuíam registros de mortes no trânsito entre 2018 e 2020. Do total de
municípios existentes no Brasil em 2021 (5.570), 5.044 apresentaram
registros de óbitos por sinistros viários nesse período. No entanto,
para assegurar a robustez da análise, foram excluídos os municípios
considerados outliers, restando um total de 4.473 municípios com dados
considerados confiáveis. As metas de redução foram estabelecidas
exclusivamente para esses municípios não outliers e por isso, neste
estudo foi realizada a comparação apenas destes municípios.

A principal base utilizada para atualização dos dados foi a
\texttt{roadtrafficdeaths}, em que \textbf{cada linha representa uma
morte por sinistro de trânsito}, contendo informações detalhadas como
data, localização e características da ocorrência. A partir dessa base,
foi realizado um \textbf{processo de agregação por município},
contabilizando o total de óbitos registrados no ano de 2023.

Em seguida, os dados agregados foram integrados à base de metas do
estudo original (realizado entre 2020 e 2021). A comparação da redução
de mortes no trânsito em 2023 em relação às metas de redução foram
calculadas da seguinte forma para cada município:

\[
\text{Redução} = \frac{\text{Número de mortes (2023)} - \text{Média de mortes (2018–2020)}}{\text{Média de mortes (2018–2020)}}
\]

\[
\text{Percentual da meta atingida} = \left( \frac{\text{Redução}}{\text{Meta estabelecida}} \right) \times 100
\]

A análise foi realizada utilizando a linguagem R, e os resultados foram
sistematizados em um dashboard desenvolvido com o pacote \texttt{Shiny}.
A ferramenta permite a visualização interativa dos resultados por
município e por estado, incluindo gráficos comparativos e tabelas, com o
intuito de facilitar a visualização do cumprimento das metas.

\bookmarksetup{startatroot}

\chapter{Resultados}\label{resultados}

\section{Resumo}\label{resumo-1}

Com o objetivo de tornar os resultados deste estudo acessíveis de forma
interativa e transparente, foi desenvolvido um \textbf{dashboard em R
utilizando o pacote Shiny}, disponível no seguinte endereço:

\url{https://beatrizmarques.shinyapps.io/cumprimento_metas/}

O painel permite a visualização detalhada do \textbf{cumprimento das
metas de redução de mortes no trânsito} estabelecidas pelo PNATRANS para
os municípios brasileiros classificados como não outliers (n = 4.473).

A interface é organizada por \textbf{unidades da federação}, e cada aba
apresenta:

\begin{itemize}
\item
  \textbf{Gráfico de dispersão}: mostra a variação proporcional do
  número de mortes em 2023 em relação à média de 2018--2020, destacando
  os municípios que atingiram ou não suas metas;
\item
  \textbf{Gráfico de barras}: exibe o percentual de municípios que
  cumpriram as metas em cada estado e capital;
\item
  ️ \textbf{Tabela interativa}: apresenta, para cada município,
  informações como população, frota, número de mortes em 2023, meta
  individualizada e se foi atingida;
\end{itemize}

O dashboard foi desenvolvido para ser uma \textbf{ferramenta de
visualização do cumprimento das metas}, possibilitando a identificação
de áreas críticas e de boas práticas locais.

\section{Abas}\label{abas}

O dashboard foi estruturado com um menu lateral que organiza o conteúdo
por \textbf{unidades da federação}, facilitando a navegação e a análise
regionalizada do cumprimento das metas.

Cada aba representa um estado brasileiro (ou o consolidado nacional
``Brasil'') e, ao ser selecionada, exibe um conjunto de visualizações e
informações específicas para os municípios daquela unidade federativa.
Essa estrutura permite ao usuário realizar consultas direcionadas,
comparações regionais e identificar padrões de desempenho na redução de
mortes no trânsito.

\section{Conteúdo}\label{conteuxfado}

Para cada aba (UF), o conteúdo apresentado inclui:

\begin{itemize}
\item
  Média de mortes no trânsito (2018-2020)
\item
  Número de mortes no trânsito (2023);
\item
  Percentual de municípios que já atingiram a meta;
\item
  Total de municípios contabilizados no cálculo (não outliers);
\item
  Qual era a meta estabelecida para o total;
\item
  Percentual de redução ou aumento no número de mortes;
\item
  Gráfico de barras por estado;
\item
  Gráfico de barras por capital;
\item
  Gráfico de dispersão: meta estabelecida x meta atingida;
\item
  Tabela interativa.
\end{itemize}

\subsection{Gráfico de Barras}\label{gruxe1fico-de-barras}

\bookmarksetup{startatroot}

\chapter*{References}\label{references}
\addcontentsline{toc}{chapter}{References}

\markboth{References}{References}

\phantomsection\label{refs}
\begin{CSLReferences}{0}{1}
\end{CSLReferences}




\end{document}
